% DO NOT MOVE KEYWORDS BELOW ABSTRACT BLOCK!
% Otherwise, LaTeX will not render them (possible a bug in the template class file)
\keyword{Digital Twin}
\keyword{AI Agent Protocol}
\keyword{Semantic Interoperability}
\keyword{Digital Ecosystem}
\keyword{Cross-domain Collaboration}
\keyword{Smart City}
\keyword{Energy Grid}

\begin{abstract}

Digital Twins (DTs) integrate real-time data streams, simulation models, and Artificial Intelligence (AI)-driven analytics across domains such as
    smart cities,
    energy grids,
    healthcare, and
    manufacturing, 
    enabling monitoring, predictive analytics, and optimization.
As DT adoption grows, cross-domain collaboration creates significant interoperability challenges, including
    technical heterogeneity,
    semantic incompatibility between domain-specific ontologies, and
    high integration maintenance burdens.

While traditional approaches rely on standardized protocols and unified data models,
    this research proposes agent-mediated semantic interoperability.
By leveraging AI Agent Protocols within Digital Ecosystems,
    autonomous DT systems can maintain domain-specific optimizations
    while dynamically negotiating semantic mappings for cross-domain collaboration,
    preserving operational integrity without disruptive integration changes.

We propose AgenTwin, an AI Agent Framework for Cooperative Digital Twin Interoperability, comprising
    a novel Digital Ecosystem (DE) Architecture and
    a proof-of-concept implementation over open-source DT frameworks.
Validation through realistic Smart City and Energy Grid cross-domain scenarios demonstrates
    the viability of solving complex optimization challenges that neither domain could address in isolation.
\end{abstract}

\translatedtitle{AgenTwin: Uma Estrutura de Agentes de IA para Interoperabilidade Cooperativa de Gêmeos Digitais}
\translatedkeyword{Gêmeo Digital}
\translatedkeyword{Protocolo de Agente de IA}
\translatedkeyword{Interoperabilidade Semântica}
\translatedkeyword{Ecossistema Digital}
\translatedkeyword{Colaboração Interdomínios}
\translatedkeyword{Cidade Inteligente}
\translatedkeyword{Rede de Energia}

\begin{translatedabstract}

Gêmeos Digitais (GDs) integram fluxos de dados em tempo real, modelos de simulação e análises orientadas por Inteligência Artificial (IA) em domínios como
    cidades inteligentes,
    redes de energia,
    saúde, e
    manufatura,
    viabilizando monitoramento, análises preditivas e otimização.
Com o crescimento da adoção de GDs, a colaboração entre domínios cria desafios significativos de interoperabilidade, incluindo
    heterogeneidade técnica,
    incompatibilidade semântica entre ontologias específicas de domínio, e
    altos custos de manutenção de integração.

Enquanto as abordagens tradicionais dependem de protocolos padronizados e modelos de dados unificados, esta pesquisa propõe a interoperabilidade semântica mediada por agentes.
Ao aproveitar Protocolos de Agentes de IA em Ecossistemas Digitais, sistemas de GDs autônomos mantêm otimizações específicas de domínio
    enquanto negociam dinamicamente mapeamentos semânticos para colaboração entre domínios,
    preservando a integridade operacional sem mudanças disruptivas de integração.

Propomos o AgenTwin, uma Estrutura de Agentes de IA para Interoperabilidade Cooperativa de Gêmeos Digitais, composta por
    uma nova Arquitetura de Ecossistema Digital (ED) e
    uma implementação de prova de conceito sobre frameworks de GDs de código aberto.
A validação por meio de cenários realistas de Cidade Inteligente e Rede de Energia demonstra
    a viabilidade de resolver desafios complexos de otimização que nenhum domínio poderia abordar de forma isolada.
\end{translatedabstract}
